\chapter{Resultados}\label{cap:resultados}

Neste capítulo, apresentaram-se os produtos da execução espacial com extração das linhas 18 a 71
dos panoramas, conforme o \autoref{cap:metodologia}. Foram examinados os sinais, os campos
wavelet e duas medidas complementares: o coeficiente de Pearson $r$ entre descritores
espectrais individuais e o escore de coerência $S_{xy}^{(0{,}05)}$ entre pares de sinais.
As figuras e tabelas foram obtidas dos artefatos ativos da execução
\codigo{v16\_spatial2\_18\_71}; os resultados arquivados não integraram esta apresentação.

% =========================================================
\section{Conjunto processado}\label{sec:res-visao}
% =========================================================

Foram analisados 26 ensaios: dez de \textit{churn}, nove de anular ondulado, três de
\textit{plug}, três de \textit{slug} e um de bolhas dispersas. Obtiveram-se 325 pares
distintos, divididos em 87 comparações intrapadrão e 238 comparações entre padrões.
A \autoref{tab:res-ensaios} identificou as aquisições na mesma ordem das matrizes
apresentadas adiante. O sufixo comum \codigo{\_001} foi omitido apenas para facilitar
a leitura. Os nomes \textit{wavy annular} e \textit{dispersed bubble}, preservados nas
figuras originais, designaram anular ondulado e bolhas dispersas.

\begin{table}[htbp]
  \centering
  \caption{Ensaios incluídos e extensão dos panoramas completos.}
  \label{tab:res-ensaios}
  \small
  \begin{tabular}{r l l r}
    \hline
    \textbf{Nº} & \textbf{Ensaio} & \textbf{Padrão} & \textbf{Extensão (px)}\\
    \hline
    1 & m03q05 & \textit{Churn} & 9586\\
    2 & m04q10 & \textit{Churn} & 8829\\
    3 & m04q20 & \textit{Churn} & 9344\\
    4 & m05q10 & \textit{Churn} & 7001\\
    5 & m05q20 & \textit{Churn} & 7391\\
    6 & m05q40 & \textit{Churn} & 9937\\
    7 & m06q10 & \textit{Churn} & 7584\\
    8 & m06q20 & \textit{Churn} & 10044\\
    9 & m07q20 & \textit{Churn} & 8255\\
    10 & m07q40 & \textit{Churn} & 10349\\
    \hline
    11 & m03q10 & Anular ondulado & 14224\\
    12 & m03q20 & Anular ondulado & 23079\\
    13 & m03q40 & Anular ondulado & 20743\\
    14 & m04q40 & Anular ondulado & 14960\\
    15 & m04q80 & Anular ondulado & 22636\\
    16 & m05q80 & Anular ondulado & 17524\\
    17 & m06q40 & Anular ondulado & 14772\\
    18 & m06q80 & Anular ondulado & 18665\\
    19 & m07q80 & Anular ondulado & 16658\\
    \hline
    20 & m05q01 & \textit{Plug} & 5940\\
    21 & m06q01 & \textit{Plug} & 7798\\
    22 & m07q05 & \textit{Plug} & 7358\\
    23 & m04q01 & \textit{Slug} & 5071\\
    24 & m05q05 & \textit{Slug} & 6323\\
    25 & m06q05 & \textit{Slug} & 7787\\
    26 & m07q01 & Bolhas dispersas & 3171\\
    \hline
  \end{tabular}
  \fonte{Elaborado pelo autor a partir do índice de ensaios da execução.}
\end{table}

Cada ensaio originou um panorama com seu sinal e um descritor espectral fixo.
Cada par originou um conjunto de campos wavelet sobre o prefixo comum dos dois sinais,
isto é, sobre a extensão do menor deles. A diferença entre essas duas entradas foi
necessária para interpretar os resultados: o coeficiente de Pearson $r$ utilizou
panoramas completos; o escore de coerência $S_{xy}^{(0{,}05)}$ utilizou os prefixos comuns.

% =========================================================
\section{Panoramas e sinais espaciais}\label{sec:res-panorama}\label{sec:res-sinais}
% =========================================================

A \autoref{fig:res-panorama-plug} apresentou a transformação da imagem em sinal:
o painel superior mostrou o panorama completo e a faixa de extração; o central apresentou
a intensidade média por coluna; o inferior mostrou o sinal após inversão de intensidade,
remoção da tendência linear e padronização. No exemplo de \textit{plug}, as regiões
alongadas da imagem corresponderam a trechos extensos do sinal, separados por mudanças
de intensidade. Na \autoref{fig:res-panorama-bolhas}, de bolhas dispersas, observaram-se
variações mais localizadas, que apareceram como picos após o pré-processamento.

\begin{figure}[p]
  \centering
  \includegraphics[width=\textwidth,height=0.76\textheight,keepaspectratio]{figuras/v16/plug_m05q01_001.png}
  \caption{Panorama completo e sinais de \textit{plug}, ensaio m05q01.
  As linhas horizontais delimitaram a faixa de extração 18--71.}
  \label{fig:res-panorama-plug}
  \fonte{Elaborado pelo autor, execução espacial 18--71.}
\end{figure}

\begin{figure}[p]
  \centering
  \includegraphics[width=\textwidth,height=0.76\textheight,keepaspectratio]{figuras/v16/dispersed_bubble_m07q01_001.png}
  \caption{Panorama completo e sinais de bolhas dispersas, ensaio m07q01.
  A aquisição participou de 25 pares entre padrões e não teve par intrapadrão.}
  \label{fig:res-panorama-bolhas}
  \fonte{Elaborado pelo autor, execução espacial 18--71.}
\end{figure}

Esses painéis permitiram relacionar as variações do sinal às imagens que lhes deram
origem. A extensão horizontal foi expressa em pixels do panorama reconstruído;
não representou uma medição calibrada do tamanho das bolhas. A padronização conservou
variações relativas de intensidade, sem transformar o brilho em fração de vazio.

% =========================================================
\section{Campos espectrais e exemplos intrapadrão}\label{sec:res-wps}\label{sec:res-wco}
% =========================================================

As \autoref{fig:res-wps}, \autoref{fig:res-wps-annular},
\autoref{fig:res-wps-plug} e \autoref{fig:res-wps-slug} apresentaram um exemplo
intrapadrão de cada grupo com pelo menos duas aquisições. Cada figura reuniu os panoramas utilizados,
os sinais normalizados, os espectros de potência wavelet (WPS), o módulo do espectro
cruzado wavelet (WCS) e a coerência wavelet (WCO).

Nos mapas, o eixo horizontal representou o atraso da correlação, indicado por \textit{lag}
nas figuras; o eixo vertical representou o comprimento de onda $\lambda$, em pixels.
Diferentemente da extensão do panorama, o atraso indicou quanto um sinal foi
deslocado em relação ao outro durante a correlação. As curvas laterais das WPS
resumiram a potência média sobre todos os atrasos. Elas não foram os descritores usados
no coeficiente de Pearson $r$, que foram extraídos da WPS no atraso nulo e
transformados conforme a \autoref{eq:met-residuos}.

\begin{table}[htbp]
  \centering
  \caption{Métricas dos pares ilustrados. O coeficiente de Pearson $r$ utilizou
  panoramas completos; o escore de coerência utilizou o prefixo comum.}
  \label{tab:res-exemplos}
  \small
  \setlength{\tabcolsep}{5pt}
  \begin{tabular}{l l r r r}
    \hline
    \textbf{Padrão} & \textbf{Par} & \textbf{Prefixo (px)} &
    \shortstack{\textbf{Coeficiente}\\\textbf{de Pearson $r$}} &
    \shortstack{\textbf{Escore de coerência}\\$S_{xy}^{(0{,}05)}$}\\
    \hline
    \textit{Churn} & m04q20 / m07q40 & 9344 & 0,9693 & 0,1799\\
    Anular ondulado & m06q40 / m07q80 & 14772 & 0,9870 & 0,1662\\
    \textit{Plug} & m05q01 / m06q01 & 5940 & 0,9670 & 0,2845\\
    \textit{Slug} & m04q01 / m05q05 & 5071 & 0,9355 & 0,4821\\
    \hline
  \end{tabular}
  \fonte{Elaborado pelo autor a partir da tabela de pares da execução.}
\end{table}

O exemplo anular apresentou coeficiente de Pearson $r=0{,}9870$ e escore de coerência
$S_{xy}^{(0{,}05)}=0{,}1662$. No exemplo de \textit{slug}, o coeficiente de Pearson
$r=0{,}9355$ permaneceu elevado, enquanto o escore de coerência
$S_{xy}^{(0{,}05)}=0{,}4821$ foi maior. Os mapas mostraram, portanto, que valores
semelhantes do coeficiente de Pearson $r$ puderam acompanhar campos de coerência
distintos. As médias de todos os pares, apresentadas a seguir, permitiram avaliar
se esse comportamento se manteve além dos exemplos.

\begin{figure}[p]
  \centering
  \includegraphics[width=\textwidth,height=0.86\textheight,keepaspectratio]{figuras/v16/churn_m04q20_001_m07q40_001.png}
  \caption{Par intrapadrão de \textit{churn}, m04q20 e m07q40: panoramas, sinais e campos WPS, WCS e WCO.}
  \label{fig:res-wps}
  \fonte{Elaborado pelo autor, execução espacial 18--71.}
\end{figure}

\begin{figure}[p]
  \centering
  \includegraphics[width=\textwidth,height=0.86\textheight,keepaspectratio]{figuras/v16/wavy_annular_m06q40_001_m07q80_001.png}
  \caption{Par intrapadrão de anular ondulado, m06q40 e m07q80. Mesma disposição da \autoref{fig:res-wps}.}
  \label{fig:res-wps-annular}
  \fonte{Elaborado pelo autor, execução espacial 18--71.}
\end{figure}

\begin{figure}[p]
  \centering
  \includegraphics[width=\textwidth,height=0.86\textheight,keepaspectratio]{figuras/v16/plug_m05q01_001_m06q01_001.png}
  \caption{Par intrapadrão de \textit{plug}, m05q01 e m06q01. Mesma disposição da \autoref{fig:res-wps}.}
  \label{fig:res-wps-plug}
  \fonte{Elaborado pelo autor, execução espacial 18--71.}
\end{figure}

\begin{figure}[p]
  \centering
  \includegraphics[width=\textwidth,height=0.86\textheight,keepaspectratio]{figuras/v16/slug_m04q01_001_m05q05_001.png}
  \caption{Par intrapadrão de \textit{slug}, m04q01 e m05q05. Mesma disposição da \autoref{fig:res-wps}.}
  \label{fig:res-wps-slug}
  \fonte{Elaborado pelo autor, execução espacial 18--71.}
\end{figure}

\clearpage
% =========================================================
\section{Perfis individuais e semelhança intrapadrão}\label{sec:res-intra}
% =========================================================

A \autoref{fig:res-perfis} apresentou, à esquerda, as WPS no atraso nulo dos
panoramas completos, divididas pelo máximo apenas para exibição. À direita apareceram
os resíduos espectrais efetivamente usados no coeficiente de Pearson $r$, após
remoção da reta ajustada ao logaritmo da WPS em função do logaritmo do comprimento
de onda $\lambda$. Os resíduos foram centralizados e divididos pela sua norma
euclidiana, isto é, pela raiz da soma de seus quadrados, formando a versão unitária
exibida. Essa normalização não alterou o coeficiente de Pearson $r$.
As linhas finas identificaram a variação entre ensaios; a linha
espessa correspondeu à mediana em cada comprimento de onda $\lambda$.

Os três descritores de \textit{plug} apresentaram formas próximas; em \textit{slug},
uma das curvas se afastou das demais. No anular ondulado, a concentração de várias
curvas foi acompanhada de um perfil deslocado para comprimentos de onda menores.
Esse registro, m03q40, foi mantido e examinado na \autoref{sec:disc-qualidade}.
A \autoref{tab:res-intra} quantificou essas observações sem reduzir cada padrão a
um único exemplo.

\begin{figure}[p]
  \centering
  \includegraphics[width=\textwidth,height=0.83\textheight,keepaspectratio]{figuras/v16/perfis_comparados.png}
  \caption{Perfis de WPS no atraso nulo e resíduos de norma unitária para o coeficiente de Pearson $r$.
  Linhas finas: ensaios; linhas espessas: medianas.}
  \label{fig:res-perfis}
  \fonte{Elaborado pelo autor, execução espacial 18--71.}
\end{figure}

\begin{table}[htbp]
  \centering
  \caption{Médias intrapadrão das duas medidas.}
  \label{tab:res-intra}
  \begin{tabular}{l r r r r}
    \hline
    \textbf{Padrão} & \textbf{Ensaios} & \textbf{Pares} &
    \shortstack{\textbf{Coeficiente de}\\\textbf{Pearson $r$ médio}} &
    \shortstack{\textbf{Escore de coerência}\\\textbf{$S_{xy}^{(0{,}05)}$ médio}}\\
    \hline
    \textit{Churn} & 10 & 45 & 0,8447 & 0,2196\\
    Anular ondulado & 9 & 36 & 0,9045 & 0,1247\\
    \textit{Plug} & 3 & 3 & 0,9618 & 0,2836\\
    \textit{Slug} & 3 & 3 & 0,7354 & 0,3681\\
    Bolhas dispersas & 1 & 0 & --- & ---\\
    \hline
  \end{tabular}
  \fonte{Elaborado pelo autor a partir do resumo por padrão da execução.}
\end{table}

O maior coeficiente de Pearson $r$ médio intrapadrão ocorreu em \textit{plug}
($0{,}9618$), seguido de anular ondulado ($0{,}9045$), \textit{churn}
($0{,}8447$) e \textit{slug} ($0{,}7354$). O escore de coerência
$S_{xy}^{(0{,}05)}$ médio apresentou outra ordem, com o maior valor em
\textit{slug} e o menor em anular ondulado. Para bolhas dispersas, o único
ensaio não permitiu estimar variabilidade intrapadrão.

\clearpage
% =========================================================
\section{Comparações entre padrões}\label{sec:res-inter}
% =========================================================

A \autoref{fig:res-pearson} reuniu os coeficientes de Pearson $r$ dos 26 ensaios.
Os blocos sobre a diagonal corresponderam aos pares do mesmo padrão; os demais
blocos, aos pares entre padrões. A diagonal unitária comparou cada descritor
consigo mesmo e foi excluída das médias. Houve coeficientes de Pearson $r$
elevados em ambos os tipos de bloco. A linha e a coluna de m03q40 e as de
m07q01 apresentaram valores menores em diversas comparações.

\begin{figure}[p]
  \centering
  \includegraphics[width=\textwidth,height=0.80\textheight,keepaspectratio]{figuras/v16/matriz_pearson_todos.png}
  \caption{Matriz do coeficiente de Pearson $r$ entre os descritores dos panoramas completos.
  As linhas brancas separaram os padrões na ordem da \autoref{tab:res-ensaios}.}
  \label{fig:res-pearson}
  \fonte{Elaborado pelo autor, execução espacial 18--71.}
\end{figure}

A \autoref{tab:res-inter} detalhou as dez combinações entre padrões. O coeficiente
de Pearson $r$ médio entre \textit{churn} e \textit{plug} foi $0{,}8988$,
superior à média intrapadrão de \textit{churn}, $0{,}8447$. As comparações com
bolhas dispersas apresentaram médias menores do coeficiente de Pearson $r$,
mas envolveram sempre a mesma aquisição desse padrão. Assim, essa diferença
descreveu um registro particular, sem estimar a variabilidade de bolhas dispersas.

\begin{table}[htbp]
  \centering
  \caption{Médias das comparações entre padrões.}
  \label{tab:res-inter}
  \small
  \begin{tabular}{l r r r}
    \hline
    \textbf{Combinação} & \textbf{Pares} &
    \shortstack{\textbf{Coeficiente de}\\\textbf{Pearson $r$ médio}} &
    \shortstack{\textbf{Escore de coerência}\\\textbf{$S_{xy}^{(0{,}05)}$ médio}}\\
    \hline
    \textit{Churn} / anular ondulado & 90 & 0,8385 & 0,2155\\
    \textit{Churn} / \textit{plug} & 30 & 0,8988 & 0,2670\\
    \textit{Churn} / \textit{slug} & 30 & 0,8271 & 0,2980\\
    \textit{Churn} / bolhas dispersas & 10 & 0,5293 & 0,3703\\
    Anular ondulado / \textit{plug} & 27 & 0,8361 & 0,2565\\
    Anular ondulado / \textit{slug} & 27 & 0,7934 & 0,2834\\
    Anular ondulado / bolhas dispersas & 9 & 0,3666 & 0,3154\\
    \textit{Plug} / \textit{slug} & 9 & 0,8765 & 0,3121\\
    \textit{Plug} / bolhas dispersas & 3 & 0,6330 & 0,3549\\
    \textit{Slug} / bolhas dispersas & 3 & 0,5634 & 0,4601\\
    \hline
  \end{tabular}
  \fonte{Elaborado pelo autor a partir do resumo por padrão da execução.}
\end{table}

A \autoref{fig:res-escore} apresentou os escores de coerência
$S_{xy}^{(0{,}05)}$. O bloco intrapadrão anular foi predominantemente escuro,
enquanto vários blocos entre padrões apresentaram valores maiores. A diagonal
cinza foi deixada sem escore porque a estatística considerou somente aquisições
distintas. A cor representou o escore ponderado dos máximos por escala, não o
máximo isolado de cada mapa nem uma probabilidade de os dois ensaios pertencerem
ao mesmo padrão.

\begin{figure}[p]
  \centering
  \includegraphics[width=\textwidth,height=0.80\textheight,keepaspectratio]{figuras/v16/matriz_escore_wco.png}
  \caption{Matriz do escore de coerência $S_{xy}^{(0{,}05)}$ calculado nos prefixos comuns.
  A diagonal cinza não integrou as médias.}
  \label{fig:res-escore}
  \fonte{Elaborado pelo autor, execução espacial 18--71.}
\end{figure}

\clearpage
% =========================================================
\section{Contrastes globais e dispersão}\label{sec:res-limite-estat}
% =========================================================

A \autoref{tab:res-global} comparou as médias intrapadrão e entre padrões nas duas
ponderações definidas na \autoref{sec:met-agregacao}. Na média por par, cada
comparação teve peso igual. Na média balanceada, cada um dos quatro padrões
com pares intrapadrão e cada uma das dez combinações entre padrões recebeu peso
igual dentro de seu conjunto. O contraste descritivo $\Delta_h$, definido na
\autoref{eq:met-contraste}, correspondeu à média intrapadrão menos a média entre
padrões. Um contraste descritivo $\Delta_h$ positivo indicou maior valor da medida
dentro dos padrões. Os valores foram arredondados a quatro casas decimais somente
para apresentação; os contrastes descritivos $\Delta_h$ foram calculados antes do
arredondamento.

\begin{table}[htbp]
  \centering
  \caption{Médias e contrastes globais das duas medidas.}
  \label{tab:res-global}
  \small
  \begin{tabular}{l l r r r}
    \hline
    \textbf{Medida} & \textbf{Ponderação} & \textbf{Intra} & \textbf{Entre} &
    \shortstack{\textbf{Contraste}\\\textbf{descritivo $\Delta_h$}}\\
    \hline
    Coeficiente de Pearson $r$ & Por par & 0,8697 & 0,8038 & 0,0659\\
    Coeficiente de Pearson $r$ & Balanceada & 0,8616 & 0,7163 & 0,1453\\
    Escore de coerência $S_{xy}^{(0{,}05)}$ & Por par & 0,1877 & 0,2635 & $-0{,}0759$\\
    Escore de coerência $S_{xy}^{(0{,}05)}$ & Balanceada & 0,2490 & 0,3133 & $-0{,}0643$\\
    \hline
  \end{tabular}
  \fonte{Elaborado pelo autor a partir do resumo intrapadrão e entre padrões da execução.}
\end{table}

O contraste descritivo $\Delta_h$ do coeficiente de Pearson $r$ foi positivo nas duas
ponderações; o contraste descritivo $\Delta_h$ do escore de coerência
$S_{xy}^{(0{,}05)}$ foi negativo. A \autoref{fig:res-distribuicoes} mostrou a
dispersão que acompanhou essas médias. Os coeficientes de Pearson $r$ variaram
de $0{,}5360$ a $0{,}9870$ dentro dos padrões e de $0{,}0018$ a $0{,}9891$
entre padrões. Os escores de coerência $S_{xy}^{(0{,}05)}$ variaram de
$0{,}0714$ a $0{,}4821$ dentro dos padrões e de $0{,}1220$ a $0{,}5554$
entre padrões.

\begin{figure}[p]
  \centering
  \includegraphics[width=\textwidth,height=0.80\textheight,keepaspectratio]{figuras/v16/correlacoes_dentro_entre.png}
  \caption{Distribuições do coeficiente de Pearson $r$ e do escore de coerência
  $S_{xy}^{(0{,}05)}$. Cada ponto representou um par; o traço horizontal marcou a
  média e a barra vertical delimitou os 50\% centrais dos valores.}
  \label{fig:res-distribuicoes}
  \fonte{Elaborado pelo autor, execução espacial 18--71.}
\end{figure}

As distribuições apresentaram ampla sobreposição. Os 325 pares compartilharam
as mesmas 26 aquisições e, portanto, não constituíram 325 observações independentes.
Os resultados foram descritivos, sem valores de significância estatística ou
estimativa de acurácia de classificação. A interpretação desse contraste positivo
do coeficiente de Pearson $r$ e a sensibilidade às escolhas de processamento e
à qualidade das entradas foram desenvolvidas no \autoref{cap:discussao}.
